% Shared preamble for the English summaries in this repository.
\documentclass[11pt,a4paper]{article}
\usepackage[margin=2.2cm]{geometry}
\usepackage[T1]{fontenc}
\usepackage[utf8]{inputenc}
\usepackage{lmodern}
\usepackage{amsmath,amssymb}
\usepackage{graphicx}
\usepackage{booktabs}
\usepackage{enumitem}
\usepackage[hidelinks]{hyperref}
\setlist{itemsep=2pt,topsep=4pt}
\setlength{\parindent}{0pt}
\setlength{\parskip}{5pt}
\graphicspath{{figures/}}
\newcommand{\summaryheader}[3]{%
  \begin{center}
    {\Large\bfseries #2}\\[3pt]
    {\large #1}\\[3pt]
    {\small Arshia Goshtasbi \textbullet{} MSc Electrical Engineering (Control), Iran University of Science and Technology \textbullet{} #3}
  \end{center}
  \vspace{2pt}\hrule\vspace{8pt}
  \noindent\textit{English summary of my original Persian report, which is kept as a PDF in this folder.}\par}

\begin{document}
\summaryheader{Convex Optimization, Take-home Quiz 3}{Convex and Quasi-convex Functions}{Spring 1404 (2025)}

\noindent\textit{The original is a scanned handwritten answer sheet in Persian (\texttt{report\_fa.pdf}).}

The quiz had nine questions on convex, concave and quasi-convex functions. My answers:
\begin{enumerate}
  \item \textbf{Sublevel sets.} A function is quasi-convex exactly when all its sublevel sets are convex; I proved both directions. Example: $f(x)=-x^2$ on $x\ge 0$ is not convex, but its sublevel sets are intervals.
  \item \textbf{Rank of PSD matrices.} $\operatorname{rank}(X)$ on $\mathbf{S}^n_+$ is not quasi-convex. The counterexample is $X=\operatorname{diag}(1,0)$, $Y=\operatorname{diag}(0,1)$, $\theta=\tfrac12$. It is quasi-concave, which I proved through $\ker(X+Y)=\ker X\cap\ker Y$ for $X,Y\succeq 0$.
  \item \textbf{A convex function and an inequality.} $x/(1-x)$ is convex on $[0,1)$, since $f''=2/(1-x)^3>0$. Applying Jensen's inequality to it gives $\frac{a}{b+c}+\frac{b}{a+c}+\frac{c}{a+b}\ge\frac32$.
  \item \textbf{Three functions.}
    \begin{itemize}
      \item The logistic function $1/(1+e^{-2x})$ is convex for $x<0$ and concave for $x>0$.
      \item $n\operatorname{tr}(X)\log\operatorname{tr}(X)-\operatorname{tr}(X)\log\det X$ is convex on $\mathbf{S}^n_{++}$. The proof works with $t=\operatorname{tr}X$ and a perspective.
      \item $\binom{x}{r}$ is convex on $(r-1,\infty)$. I showed this through $\log\binom{x}{r}=\sum_k \log(x-k)-\log r!$.
    \end{itemize}
  \item \textbf{Quasi-concavity on an interval.} A continuous $f$ on $[a,b]$ is quasi-concave if $f(\theta a+(1-\theta)b)\ge\min\{f(a),f(b)\}$.
  \item \textbf{Separating functions.} A convex $f$ and a concave $g$ with $g\le f$ can be separated by an affine function $h$ with $g\le h\le f$.
  \item \textbf{Determinant root.} $(\det X)^{1/n}$ is concave on $\mathbf{S}^n_{++}$. I restricted it to a line $Z+tV$, which reduces it to a geometric mean of $1+t\lambda_i$.
  \item \textbf{A perspective.} $\|x\|_p^p/t^{p-1}$ ($p>1$) is convex as the perspective of the convex function $\|x\|_p^p$.
  \item \textbf{Square of a function.} $f^2$ is convex when $f$ is convex and non-negative, or concave and non-positive. The proof uses $(f^2)''=2f'^2+2ff''$.
\end{enumerate}
\end{document}
